% META: {"id": "TSPE-TRIGO-ME-002", "chapitre": "TSPE-TRIGONOMETRIE", "type_objet": "methode", "capacites_codes": ["C2"], "methodes": ["M2"], "sources_inspiration": [], "status": "approved"}
\begin{fichemethode}{M2}{Étudier une fonction trigonométrique pour optimiser}

\quandUtiliser{%
  Utiliser cette méthode lorsqu'un probleme (souvent géométrique) se ramene a l'étude des variations d'une fonction construite avec $\sin$ ou $\cos$.
}

\pasApas{%
  \begin{enumerate}
    \item \textbf{Modeliser} : exprimer la quantité a optimiser (aire, longueur, volume...) en fonction d'une variable $x$ (souvent un angle), en precisant l'intervalle d'étude.
    \item \textbf{Dériver} en utilisant $\sin'=\cos$, $\cos'=-\sin$ et, si besoin, la dérivée d'une composee $(v \circ u)' = u' \times (v' \circ u)$.
    \item \textbf{Étudier le signe} de la dérivée sur l'intervalle d'étude.
    \item \textbf{Dresser le tableau de variations}, identifier le maximum ou le minimum et la valeur de $x$ en laquelle il est atteint.
    \item \textbf{Conclure} en repondant a la question posee dans le contexte du probleme (avec les bonnes unités).
  \end{enumerate}
}

\end{fichemethode}
